\section{优化：在证明义务约束下改写程序}
\label{sec:optimization}

“\texttt{-O2} 让程序更快”只描述了意图，没有解释机制。中端优化更准确的模型是：分析先建立某个事实，变换在该事实允许的范围内改写 IR，并声明哪些分析仍然有效；后续 pass 再利用新形状继续推理。LLVM 的新 Pass Manager 按模块、调用图强连通分量、函数和循环等 IR 单元组织这一过程，分析结果会被缓存，也必须在变换破坏前提时失效\cite{llvm-newpm}。因此优化不是一串互不相干的“技巧”，而是一个对表示形状高度敏感的因果链。

本节比较两份真实工件。它们均由 Ubuntu Clang 18.1.3 从固定 C 观察载体生成：基线使用 \texttt{-O0 -Xclang -disable-O0-optnone}，优化版本使用 \texttt{-O2 -fno-inline}。禁用 \texttt{optnone} 使基线 IR 可继续交给实验 pass；\texttt{-fno-inline} 则刻意保留 \texttt{clipped\_dot} 边界，方便观察。两份文件的目标是 x86-64，而第~\ref{sec:backend} 节另行分析同一源程序面向 RV64GC 的 \texttt{-O2} 汇编。这个实验设计很重要：若目标、内联策略和优化级别一起变化，差异就无法归因。

\subsection{证据分级：看到结果不等于知道是哪一个 pass}

后文采用三种陈述强度：
\begin{description}
  \item[\observed{}] 可直接在保存的 \texttt{.ll}/\texttt{.s} 工件中定位，例如 \texttt{-O2} IR 出现了 \texttt{<4 x i32>}；
  \item[\inferred{}] 根据 LLVM 已公开的 pass 职责与前后表示作出的最佳机制解释，例如局部栈槽消失符合 mem2reg/SROA 的作用，但未保存逐 pass 快照时不能声称某一具体 pass 独自完成；
  \item[反事实条件] 解释若存在写内存、别名或不安全重排会阻止什么，用于揭示合法性边界，而不是伪造一次“失败实验”。
\end{description}

默认优化流水线由 \texttt{PassBuilder}、LLVM 版本、优化级别、目标机器、LTO/PGO 和前端插入点共同决定；\texttt{-O2} 不是跨版本稳定的 pass 名单。要建立严格因果证据，应额外保存 \texttt{-debug-pass-manager} 日志、\texttt{-print-before/after} 快照和优化备注。本仓库当前保存的是端点工件，因而表~\ref{tab:o0-o2-observed} 只把端点差异列为观察，把可能机制明确写成解释。

\begin{table}[t]
  \centering
  \caption{Clang 18.1.3 的 \texttt{-O0} 与 \texttt{-O2 -fno-inline} 端点差异}
  \label{tab:o0-o2-observed}
  \scriptsize
  \begin{tabular}{p{0.27\textwidth}p{0.29\textwidth}p{0.33\textwidth}}
    \toprule
    \texttt{-O0} 观察 & \texttt{-O2} 观察 & 合理机制解释（非单-pass 归因） \\
    \midrule
    参数及 \texttt{i}/\texttt{acc}/\texttt{term} 均有 \texttt{alloca--store--load} & 内核无 \texttt{alloca}，循环状态由 \texttt{phi} 传递 & mem2reg/SROA 与后续简化将可提升对象变为 SSA \\
    截断由多个条件分支和临时栈槽表达 & 标量尾部使用 \texttt{llvm.smin.i32} 加 \texttt{select}；向量主循环使用向量比较、\texttt{smin} 和 \texttt{select} & InstCombine、SimplifyCFG 等规范化出更易识别的最小/最大模式 \\
    每轮处理一个元素 & 主循环每轮读两个 \texttt{<4 x i32>}，即 8 项；另有标量尾循环 & LoopVectorize 生成向量主体、水平归约和余数处理 \\
    \texttt{main} 有两个栈数组及两次 \texttt{llvm.memcpy} & 直接把 \texttt{@\_\_const.main.a/b} 地址传入 & 常量传播、全局化/内存优化与死对象消除共同作用 \\
    \texttt{n} 的两级 CFG & \texttt{llvm.smin.i32} 后接 \texttt{llvm.smax.i32} & 条件值化与 min/max 惯用式识别 \\
    函数属性保守 & 内核具有 \texttt{nofree norecurse nosync nounwind memory(argmem: read)} 等 & 函数属性推断基于函数体建立调用侧可用事实 \\
    \bottomrule
  \end{tabular}
\end{table}

\subsection{第一因：从内存形式到可分析的 SSA}

Clang 的 \texttt{-O0} 内核忠实保留了源变量的“可寻址”形状：
\begin{lstlisting}[language=LLVMIR,basicstyle=\ttfamily\scriptsize]
%6  = alloca ptr, align 8       ; a
%8  = alloca i32, align 4       ; n
%11 = alloca i32, align 4       ; acc
%12 = alloca i32, align 4       ; i
%13 = alloca i32, align 4       ; term
store i32 0, ptr %11, align 4
store i32 0, ptr %12, align 4
...
%15 = load i32, ptr %12, align 4
%16 = load i32, ptr %8, align 4
%17 = icmp slt i32 %15, %16
\end{lstlisting}
这并不与“LLVM IR 是 SSA”矛盾：\texttt{\%15} 等加载结果各自只定义一次，而栈对象可以被多次写入。内存提升（memory promotion，通常称 mem2reg）寻找入口块中、地址不逸出且只由普通 load/store 使用的 \texttt{alloca}，在支配边界插入必要的 \texttt{phi}，随后沿支配树重命名值\cite{llvm-cytron1991,llvm-passdocs}。聚合体标量替换（Scalar Replacement of Aggregates, SROA）还可拆分聚合对象或切片，使其中的标量部分进入同一过程。

二者的能力边界同样重要：数组地址一旦传给未知函数、对象存在重叠切片、访问为 \texttt{volatile}，或地址以优化器不能证明安全的方式逸出，就不能简单提升。\texttt{main} 中的数组语义也不是由 mem2reg 一步变成全局常量；其初始化、只读使用和调用效应需要多个分析/变换共同证明。将所有“栈消失”都归功于 mem2reg，是一种看似具体、实则错误的解释。

\observed{}，\texttt{-O2} 内核不再含任何 \texttt{alloca}，标量尾循环用两个 \texttt{phi} 携带索引与累加值：
\begin{lstlisting}[language=LLVMIR,basicstyle=\ttfamily\scriptsize]
%48 = phi i64 [ %59, %47 ], [ %45, %44 ]
%49 = phi i32 [ %58, %47 ], [ %46, %44 ]
...
%58 = add nsw i32 %57, %49
%59 = add nuw nsw i64 %48, 1
\end{lstlisting}
一旦“当前 \texttt{i} 与 \texttt{acc}”成为直接的 def--use 链，循环分析不必再反复询问某个 load 是否可能被别名 store 改写，归纳变量和归约也就从内存谜题变成图上的模式。

\subsection{第二因：规范形状让后续分析看得见}

规范化（canonicalization）并不等于为了美观重写文本。InstCombine 倾向于在不改变 CFG 的前提下把表达式变成下游约定的惯用形；SimplifyCFG 合并冗余块、折叠分支并把合适的控制依赖变为 \texttt{select}\cite{llvm-passdocs}。在案例中，\texttt{-O0} 由 \texttt{term} 栈槽和三个分支块实现上下界，而 \texttt{-O2} 标量尾部为：
\begin{lstlisting}[language=LLVMIR,basicstyle=\ttfamily\scriptsize]
%54 = mul nsw i32 %53, %51
%55 = icmp slt i32 %54, %3
%56 = tail call i32 @llvm.smin.i32(i32 %54, i32 %4)
%57 = select i1 %55, i32 %3, i32 %56
\end{lstlisting}
其语义仍是 \(\max(lo,\min(product,hi))\)，前提是调用侧给出的边界关系与原源程序一致。该无分支形状天然逐通道（lane-wise）工作，因此同一结构可提升为 \texttt{<4 x i32>} 比较、\texttt{smin} 和 \texttt{select}。换言之，先消除“临时变量必须住在栈上”与“条件赋值必须是 CFG”这两个偶然表示，向量化所需的数据并行结构才显现出来。

全局值编号（Global Value Numbering, GVN）会把表达式按值等价类处理，并在内存依赖允许时消除冗余 load；循环不变代码外提（Loop-Invariant Code Motion, LICM）尝试把循环内不变且可安全执行/移动的运算提到循环外。这两类 pass 均依赖支配、别名和副作用信息。当前端点中没有保存足以唯一证明“某条指令正由 GVN 删除”或“某条指令正由 LICM 外提”的逐 pass 证据，因此我们不做该观察性声称；它们在因果链中的角色是说明何时冗余访存和循环不变量\emph{可以}被消除。

例如，把 \texttt{a[i]} 连续读取两次，若中间只有 \texttt{memory(none)} 调用，则第二次 load 可能复用第一次的值；若中间调用缺少内存效应属性，它可能写入 \texttt{a[i]}，这个消除便不合法。同理，可能陷阱的 load 通常不能仅因“地址表达式与 \texttt{i} 无关”就无条件提到零次迭代循环之前。优化不是代数化简加速版，而是代数、控制可达性与内存模型的交集。

\subsection{别名是向量化合法性的闸门}

别名分析（alias analysis）回答两个内存位置是否可能指向同一对象或重叠区域。它通常返回 MustAlias、NoAlias、MayAlias 等保守结果；“分析不知道”必须落入可能别名，而不能乐观当作不别名\cite{llvm-aa}。此外，函数属性中的 \texttt{memory(argmem: read)}、参数 \texttt{readonly} 以及基于类型的别名分析（Type-Based Alias Analysis, TBAA）元数据共同描述内存效应。

本案例有一个容易被忽略的好性质：内核对 \texttt{a} 和 \texttt{b} 只有读，没有向二者或第三个数组写。即使两个参数完全别名，同时从相同位置读取仍不形成循环携带的读写依赖。因此 \texttt{a} 与 \texttt{b} 没有 \texttt{noalias} 也不妨碍该次向量化；\texttt{-O2} 函数属性确实只显示 \texttt{nocapture readonly}，未声称二者不别名。这比“向量化总要求 \texttt{restrict}”更精确。

反事实地，若循环改为 \texttt{a[i] = ... b[i-1] ...}，某些重叠关系会让后一迭代读到前一迭代刚写的数据；向量执行将改变顺序。编译器此时要么证明 NoAlias，要么生成运行时别名检查并保留标量版本，要么放弃向量化。属性可以解除闸门，但虚假的 \texttt{noalias} 不是性能技巧，而是错误的程序契约。

\subsection{从标量循环到向量主体与标量尾部}

\observed{}，x86-64 \texttt{-O2} 工件首先判断 \(n>0\)，把 \texttt{n} 零扩展为 64 位，并在 \(n\ge8\) 时进入向量路径。其迭代上界
\begin{lstlisting}[language=LLVMIR,basicstyle=\ttfamily\scriptsize]
%11 = and i64 %8, 2147483640
\end{lstlisting}
把非负 \texttt{n} 向下取整到 8 的倍数。向量循环每轮进行四次宽度为 4 的 load 两对，即同时处理 8 项，并维护两个独立的 \texttt{<4 x i32>} 累加器：
\begin{lstlisting}[language=LLVMIR,basicstyle=\ttfamily\scriptsize]
%22 = load <4 x i32>, ptr %20, align 4, !tbaa !5
%23 = load <4 x i32>, ptr %21, align 4, !tbaa !5
%26 = load <4 x i32>, ptr %24, align 4, !tbaa !5
%27 = load <4 x i32>, ptr %25, align 4, !tbaa !5
%28 = mul nsw <4 x i32> %26, %22
...
%36 = add <4 x i32> %34, %18
%37 = add <4 x i32> %35, %19
%38 = add nuw i64 %17, 8
\end{lstlisting}
双累加器降低了单条归约依赖链对吞吐的限制。循环结束后先逐通道相加，再以 \texttt{llvm.vector.reduce.add.v4i32} 做水平归约；若 \(n\) 不是 8 的倍数，\texttt{\%45}/\texttt{\%46} 把已经完成的索引和部分和交给标量尾循环；若 \(n<8\)，则从 0 直接进入该尾循环。\texttt{exit} 的三路 \texttt{phi} 分别覆盖零次迭代、恰好完成向量主体、以及向量主体加标量余数。这不是“把 \texttt{i32} 改成向量类型”一项改写，而是 CFG、归纳变量、归约和余数策略的联合变换。

这次成功不能外推为所有点积都会向量化。向量器首先要证明合法性，再根据目标变换信息（Target Transform Info, TTI）评估收益\cite{llvm-vectorizers}。整数加法在位向量层面可结合，但源 C 的有符号溢出是 UB；优化器仅需保持所有已定义执行。在固定案例的可达输入上，每项已截到 \([-8,12]\)，最多 8 项，总和位于 \([-64,96]\)，所以归约无溢出。若改成浮点数，IEEE-754 舍入使加法通常不满足结合律，除非快速数学标志或语言模式明确放宽语义，不能照搬此次重排。

\subsection{目标相关的盈利性：同一中端机会并非同一机器结果}

\observed{}，面向 x86-64 的优化 IR 选择了 \texttt{<4 x i32>} 向量主体；而保存的 RV64GC 汇编由
\begin{lstlisting}[language=RISCV,basicstyle=\ttfamily\scriptsize]
clang --target=riscv64-linux-gnu -march=rv64gc -mabi=lp64d \
      -S -O2 -fno-inline ...
\end{lstlisting}
生成，目标特性中没有 RISC-V V 扩展，因此内核是标量 \texttt{lw}/\texttt{mulw}/\texttt{addw} 循环。中端与后端的边界并非“中端保证向量、后端照抄”：向量器会咨询目标合法类型、寄存器宽度、指令成本和对齐代价；某个目标上的有利形状可能在另一个目标上保留标量。

主函数也提供了成本与全程序事实的例子。\observed{}，\texttt{-O2} IR 没有把内核内联，因为命令显式使用 \texttt{-fno-inline}；因此“向量化源于内联”可被排除。与此同时，两个局部数组变为直接引用只读全局模板，\texttt{n} 的夹取变为 \texttt{smin}/\texttt{smax}，而调用仍在。这种刻意保留边界的设计使读者能同时观察跨函数属性推断和函数内循环优化。

\subsection{一条可复现的因果实验路线}

若要把本节的\inferred{}提升为逐步\observed{}，应从禁用 \texttt{optnone} 的基线开始，而不是手工猜测默认管线：
\begin{lstlisting}[basicstyle=\ttfamily\scriptsize]
clang -O0 -Xclang -disable-O0-optnone -S -emit-llvm case.c -o raw.ll
opt -S -passes='mem2reg,sroa,instcombine,simplifycfg' raw.ll -o canon.ll
opt -S -passes='default<O2>' -verify-each raw.ll -o O2.ll
opt -passes='default<O2>' -debug-pass-manager -disable-output raw.ll
clang -O2 -Rpass=loop-vectorize -Rpass-missed=loop-vectorize \
      -Rpass-analysis=loop-vectorize case.c -c
\end{lstlisting}
实际执行前应以对应版本的 \texttt{opt --print-passes} 核对名称，并记录工具版本、命令、目标三元组和输入哈希。\texttt{-verify-each} 只能捕获每一步的结构破坏；它不证明 pass 在源语言 UB、属性和目标 ABI 上的所有推理都正确。进一步的翻译验证可把“变换前 IR 精化变换后 IR”交给 SMT，但仍需注意超时、未支持特性和语义版本差异\cite{llvm-alive2}。

把整条链压缩为一句话就是：\emph{内存提升让数据流显式，规范化让模式可识别，别名/效应分析开放合法性，目标成本模型决定是否兑现，最后由向量主体、归约和标量尾部共同保持语义。} 下一节继续沿这条证据链，考察已经优化的 IR 如何变成有限寄存器、真实指令与对象字节。
